%  The AAU Poster Theme.
%  2013-05-08 v. 1.1.0
%  Copyright 2013 by Jesper Kjær Nielsen <jkn@es.aau.dk>
%
%  This is free software: you can redistribute it and/or modify
%  it under the terms of the GNU General Public License as published by
%  the Free Software Foundation, either version 3 of the License, or
%  (at your option) any later version.
%
%  This is distributed in the hope that it will be useful,
%  but WITHOUT ANY WARRANTY; without even the implied warranty of
%  MERCHANTABILITY or FITNESS FOR A PARTICULAR PURPOSE.  See the
%  GNU General Public License for more details.
%
%  You can find the GNU General Public License at <http://www.gnu.org/licenses/>.
\documentclass[a1paper,portrait]{baposter}
%%%%%%%%%%%%%%%%%%%%%%%%%%%%%%%%%%%%%%%%%%%%%%%%
% Language, Encoding and Fonts
% http://en.wikibooks.org/wiki/LaTeX/Internationalization
%%%%%%%%%%%%%%%%%%%%%%%%%%%%%%%%%%%%%%%%%%%%%%%%
% Select encoding of your inputs. Depends on
% your operating system and its default input
% encoding. Typically, you should use
%   Linux  : utf8 (most modern Linux distributions)
%            latin1 
%   Windows: ansinew
%            latin1 (works in most cases)
%   Mac    : applemac
% Notice that you can manually change the input
% encoding of your files by selecting "save as"
% an select the desired input encoding. 
\usepackage[utf8]{inputenc}
% Make latex understand and use the typographic
% rules of the language used in the document.
\usepackage[english]{babel}
% Use the vector font Latin Modern which is going
% to be the default font in latex in the future.
\usepackage{helvet}
% Change the default font family from roman to sans serif
\renewcommand{\familydefault}{\sfdefault} % for text
\usepackage[helvet]{sfmath} % for math
% Choose the font encoding
\usepackage[T1]{fontenc}

%%%%%%%%%%%%%%%%%%%%%%%%%%%%%%%%%%%%%%%%%%%%%%%%
% Graphics and Tables
% http://en.wikibooks.org/wiki/LaTeX/Importing_Graphics
% http://en.wikibooks.org/wiki/LaTeX/Tables
% http://pgfplots.sourceforge.net/
%%%%%%%%%%%%%%%%%%%%%%%%%%%%%%%%%%%%%%%%%%%%%%%%
% You cannot use floats in the baposter theme.
% We therefore load the caption package which provides
% the command \captionof
% Set up how figure and table captions are displayed
\usepackage{caption}
\captionsetup{
  font=small,% set font size to footnotesize
  labelfont=bf % bold label (e.g., Figure 3.2) font
}
% Make the standard latex tables look so much better
\usepackage{array,booktabs}
% For creating beautiful plots
\usepackage{pgfplots}

%%%%%%%%%%%%%%%%%%%%%%%%%%%%%%%%%%%%%%%%%%%%%%%%
% Mathematics
% http://en.wikibooks.org/wiki/LaTeX/Mathematics
%%%%%%%%%%%%%%%%%%%%%%%%%%%%%%%%%%%%%%%%%%%%%%%%
% Defines new environments such as equation,
% align and split 
\usepackage{amsmath}
% Adds new math symbols
\usepackage{amssymb}

%%%%%%%%%%%%%%%%%%%%%%%%%%%%%%%%%%%%%%%%%%%%%%%%
% Colours
% http://en.wikibooks.org/wiki/LaTeX/Colors
%%%%%%%%%%%%%%%%%%%%%%%%%%%%%%%%%%%%%%%%%%%%%%%%
\selectcolormodel{RGB}
% define the three aau colors
\definecolor{aaublue1}{rgb}{0.01,0.55,0.55}% dark blue
\definecolor{aaublue2}{rgb}{0.005,0.27,0.27} % light blue
\definecolor{aaublue3}{RGB}{204,233,234} % lighter blue
\let\oldtextbf\textbf
\renewcommand{\textbf}[1]{\textcolor{aaublue1}{\oldtextbf{#1}}}
%%%%%%%%%%%%%%%%%%%%%%%%%%%%%%%%%%%%%%%%%%%%%%%%
% Lists
% http://en.wikibooks.org/wiki/LaTeX/List_Structures
%%%%%%%%%%%%%%%%%%%%%%%%%%%%%%%%%%%%%%%%%%%%%%%%
% Easier configuration of lists
\usepackage{enumitem}
%configure itemize
\setlist{%
  topsep=0pt,% set space before and after list
  noitemsep,% remove space between items
  labelindent=\parindent,% set the label indentation to the paragraph indentation
  leftmargin=*,% remove the left margin
  font=\color{aaublue1}\normalfont, %set the colour of all bullets, numbers and descriptions to aaublue1
}
% use set<itemize,enumerate,description> if you have an older latex distribution
\setitemize[1]{label={\raise1.25pt\hbox{$\blacktriangleright$}}}
\setitemize[2]{label={\scriptsize\raise1.25pt\hbox{$\blacktriangleright$}}}
\setitemize[3]{label={\raise1.25pt\hbox{$\star$}}}
\setitemize[4]{label={-}}
%\setenumerate[1]{label={\theenumi.}}
%\setenumerate[2]{label={(\theenumii)}}
%\setenumerate[3]{label={\theenumiii.}}
%\setenumerate[4]{label={\theenumiv.}}
%\setdescription{font=\color{aaublue1}\normalfont\bfseries}

% use setlist[<itemize,enumerate,description>,<level>] if you have a newer latex distribution
%\setlist[itemize,1]{label={\raise1.25pt\hbox{$\blacktriangleright$}}}
%\setlist[itemize,2]{label={\scriptsize\raise1.25pt\hbox{$\blacktriangleright$}}}
%\setlist[itemize,3]{label={\raise1.25pt\hbox{$\star$}}}
%\setlist[itemize,4]{label={-}}
%\setlist[enumerate,1]{label={\theenumi.}}
%\setlist[enumerate,2]{label={(\theenumii)}}
%\setlist[enumerate,3]{label={\theenumiii.}}
%\setlist[enumerate,4]{label={\theenumiv.}}
%\setlist[description]{font=\color{aaublue1}\normalfont\bfseries}

%%%%%%%%%%%%%%%%%%%%%%%%%%%%%%%%%%%%%%%%%%%%%%%%
% Misc
%%%%%%%%%%%%%%%%%%%%%%%%%%%%%%%%%%%%%%%%%%%%%%%%
% change/remove some names
\addto{\captionsenglish}{
  %remove the title of the bibliograhpy
  \renewcommand{\refname}{\vspace{-0.7em}}
  %change Figure to Fig. in figure captions
  \renewcommand{\figurename}{Fig.}
}
% create links
\usepackage{url}
%note that the hyperref package is currently incompatible with the baposter class

%%%%%%%%%%%%%%%%%%%%%%%%%%%%%%%%%%%%%%%%%%%%%%%%
% Macros
%%%%%%%%%%%%%%%%%%%%%%%%%%%%%%%%%%%%%%%%%%%%%%%%
\newcommand{\alert}[1]{{\color{aaublue1}#1}}

%%%%%%%%%%%%%%%%%%%%%%%%%%%%%%%%%%%%%%%%%%%%%%%%
% Document Start 
%%%%%%%%%%%%%%%%%%%%%%%%%%%%%%%%%%%%%%%%%%%%%%%%
\begin{document}
%%%%%%%%%%%%%%%%%%%%%%%%%%%%%%%%%%%%%%%%%%%%%%%%
% Some changes that cannot be made in the preamble
%%%%%%%%%%%%%%%%%%%%%%%%%%%%%%%%%%%%%%%%%%%%%%%%
% set the background of the poster
\background{
  \begin{tikzpicture}[remember picture,overlay]%
    %the poster background color
    \draw[fill=white,draw=aaublue3] (current page.north west) rectangle (current page.south east);
    %the header
    \fill [fill=aaublue3] (current page.north west) rectangle ([yshift=-\headerheight] current page.north east);
  \end{tikzpicture}
}
% if you want to reduce the space before and after equations, use and adjust
% the following lines
%\addtolength{\abovedisplayskip}{-2mm}
%\addtolength{\belowdisplayskip}{-2mm}

%%%%%%%%%%%%%%%%%%%%%%%%%%%%%%%%%%%%%%%%%%%%%%%%
% General poster setup
%%%%%%%%%%%%%%%%%%%%%%%%%%%%%%%%%%%%%%%%%%%%%%%%
\begin{poster}{
  %general options for the poster
  grid=false,
  columns=2,
%  colspacing=4.2mm,
  headerheight=0.075\textheight,
  background=user,
%  bgColorOne=red!42, %is used when background != user and none
%  bgColorTwo=green!42, %is used when background is shaded
  eyecatcher=true,
  %posterbox options
  headerborder=closed,
  borderColor=aaublue1,
%  headershape=roundedright,
%  headershade=plain,
  headerColorOne=white,
%  headerColortwo=yellow!42, %is used when the header background is shaded
  textborder=rectangle,
  boxshade=plain,
  boxColorOne=white,
%  boxColorTwo=cyan!42,%is used when the text background is shaded
  headerFontColor=aaublue2,
  headerfont=\large\sf\bf,
  linewidth=1pt
}
%the Eye Catcher (the logo on the left)
{
  %this can be commented out or replaced by a company/department logo
  \raisebox{-0.25\headerheight}[0pt][0pt]{%
      \makebox[.75\headerheight][c]{
  \includegraphics[height=.75\headerheight]{figs/logo.ps} }}
  
}
%the poster title
{\color{aaublue2}\huge\bf
%    \vspace*{-0.1em}
On-Chip Communication Fault-Resilient Adaptive  Architectures and Algorithms  for 3D-IC Technologies        
}
%the author(s)
{\color{aaublue2} \small
%  \vspace{.8em}
  Professor Abderazek Ben Abdallah, d8162103 Nam Khanh Dang | \textbf{ Adaptive Systems Laboratory}
  
}
%the logo (the logo on the right)
{
  %this can be commented out or replaced by a company/department logo
%  \includegraphics[height=.75\headerheight]{figs/asl.png}
  \raisebox{-0.25\headerheight}[0pt][0pt]{%
    \makebox[.75\headerheight][c]{
        \includegraphics[height=.75\headerheight]{figs/asl.pdf} }}
}

%%%%%%%%%%%%%%%%%%%%%%%%%%%%%%%%%%%%%%%%%%%%%%%%
% the actual content of the poster begins here
%%%%%%%%%%%%%%%%%%%%%%%%%%%%%%%%%%%%%%%%%%%%%%%%

\begin{posterbox}[name=intro,column=0,row=0]{Era of Many-core Chips}
Hundreds of cores integrated on chip:
\begin{itemize}
  \item Technology scaling.
  \item 3D integration.
  \item Several examples:
  \begin{itemize}
      \item STMicro P2012/STHORM 
      \item picoChip 
      \item Scale-out
      
      \item Tilera Tile GX, Tile Pro 
      \item Intel Polari 
  \end{itemize}
  \item Complex apps, stringent constraints:
  \begin{itemize}
      \item Massively parallel applications
      \item Performance, power, scalability, reliability are key challenges.
  \end{itemize}
\end{itemize}

\begin{tikzpicture}[remember picture,overlay]
\node at (10,2.8) {\includegraphics[width=5cm]{figs/intro_pic.pdf}};
\end{tikzpicture}
\end{posterbox}

\begin{posterbox}[name=3dnoc,column=0,below=intro]{3D Network-on-Chip}
    \begin{center}
        \includegraphics[width=1\linewidth]{figs/NoC2r.pdf}\\
        \textbf{\small System architecture: (a) Network; (b) Vertical connection; (c) Router; (d) Wrapped router.}
    \end{center}

\end{posterbox}



\begin{posterbox}[name=demand,column=0, below=3dnoc]{Reliability Demand}
\begin{itemize}
    \item NoCs are exposed to a variety of manufacturing and design factors making them vulnerable to different faults that cause corrupted message transfer or even catastrophic system failures.
    \item The single-point-failure nature of NoC introduces a big concern to their reliability as they are the sole communication medium. 
    \item Several works have dealt with separated types of faults; however, there is a demand of a comprehensively reliable 3D NoC system which can handle soft errors, hard faults and TSV defects.
\end{itemize}
\end{posterbox}



\begin{posterbox}[name=tech,column=1]{Technology Features}
%        \includegraphics[width=1\linewidth]{figs/router.pdf} \\
        
        \begin{center}
            \includegraphics[width=.8\linewidth]{figs/case1.pdf} \\
            \vspace*{-0.0cm}\textbf{\small TSV Sharing Process.}
        \end{center}
        
         \begin{center}
             \includegraphics[width=.85\linewidth]{figs/hft.pdf} \\
            \vspace*{-0cm}\textbf{\small Hard Fault Tolerance: (a) Random Access Buffer; (b) Bypass-Link-on-Demand.}
        \end{center}
        
        \begin{center}
            \includegraphics[width=.8\linewidth]{figs/sft.pdf} \\
            \vspace*{-0cm}\textbf{\small Soft Error Tolerance: Pipeline Computation Redundancy.}
        \end{center}
\end{posterbox}

\begin{posterbox}[name=refs,column=1,span=1,below=tech]{References}

    \bibliographystyle{plain}
    \begin{thebibliography}{1}% Simple bibliography with widest label of 1
        \itemsep=-0.01em% Save space between the separation
        \setlength{\baselineskip}{0.4em}% Save space with longer lines
        \bibitem{baposter} Dang Nam Khanh: \emph{\color{aaublue2}Development of On-Chip Communication Fault-Resilient Adaptive Architectures and Algorithms for 3D-IC Technologies}, Doctoral Thesis, The University of Aizu, 2017. 
        \bibitem{pgfplots} Akram Ben Ahmed: \emph{\color{aaublue2}High-throughput Architecture and Routing Algorithms Towards the Design of Reliable Mesh-based Many-Core Network-on-Chip Systems}, Doctoral Thesis, The University of Aizu, 2015. 
    \end{thebibliography}

\end{posterbox}

\begin{posterbox}[name=perf,column=0,above=bottom,span=2]{Result: Reliability, Performance and Layout }
\begin{center}
    \includegraphics[width=1\linewidth]{figs/reli.pdf} \\
    \textbf{\small Reliability evaluation. Layer size: (a) $\mathbf{16\times16}$; (b) $\mathbf{32\times32}$; (c) $\mathbf{64\times64}$. }
\end{center}
        \includegraphics[width=.75\linewidth]{figs/perf.pdf}\hspace*{0.5cm} \includegraphics[width=.22\linewidth]{figs/layout.pdf} \\
        \textbf{\small\hspace*{8cm}Performance evaluation. \hspace*{9cm} Layout of a $\mathbf{2\times2}$ layer.}
\end{posterbox}
\end{poster}
\end{document}
